\chapter{A Pesquisa em Compiladores}

\section{Introdução}

Um compilador é um dos softwares mais críticos de um ecossistema de
desenvolvimento: todo programa escrito em uma linguagem de alto nível
passa por ele. Um bug no compilador pode silenciosamente introduzir
erros em código que parece correto, comprometer a segurança de sistemas
inteiros ou impedir otimizações que o programador assume presentes.
Por isso, a pesquisa em compiladores vai muito além da construção de
compiladores corretos por construção: ela estuda como \emph{testá-los},
como \emph{prová-los formalmente corretos}, como \emph{enriquecer as
linguagens-fonte} com sistemas de tipos mais expressivos, e como
\emph{representar e transformar programas} de maneira eficiente e
verificável.

Este capítulo apresenta uma panorâmica dos tópicos de pesquisa mais
ativos na área, com o objetivo de indicar caminhos para aprofundamento.
Cada seção apresenta a motivação, os conceitos centrais e exemplos
concretos, seguidos de referências à literatura.

\section{Teste de Compiladores}

\subsection{O Problema}

Compiladores são programas grandes e complexos. O GCC possui mais de
15 milhões de linhas de código; o LLVM, mais de 20 milhões. Encontrar
bugs em compiladores (especialmente \emph{wrong-code bugs}, nos quais
o compilador gera código que produz resultado errado mas não trava)
é notoriamente difícil com testes manuais.

\subsection{Geração Aleatória de Programas (Fuzzing)}

A abordagem mais influente é gerar programas válidos aleatoriamente e
verificar se o compilador se comporta corretamente ao compilá-los. A
ferramenta \textbf{CSmith} [Yang et al., 2011] gera programas C
aleatórios com semântica definida (sem comportamento indefinido) e
descobriu centenas de bugs em GCC e LLVM. A ideia central é simples:

\begin{enumerate}
  \item Gerar um programa aleatório \(P\) com saída bem definida.
  \item Compilar \(P\) com múltiplos compiladores e com múltiplos
    níveis de otimização.
  \item Comparar as saídas: discrepâncias indicam bugs.
\end{enumerate}

Esse princípio chama-se \textbf{teste diferencial} (\emph{differential
testing}). Ele não requer um oráculo que saiba a resposta correta:
basta que dois compiladores corretos concordem.

\subsection{Teste Baseado em Propriedades}

Uma abordagem complementar é o \textbf{teste baseado em propriedades}
(\emph{property-based testing}), popularizado pela biblioteca QuickCheck
[Claessen \& Hughes, 2000]: em vez de casos de teste fixos, o
programador especifica \emph{propriedades} que o compilador deve
satisfazer, e a ferramenta gera inputs aleatórios para falsificá-las.

Por exemplo, a propriedade de \textbf{equivalência semântica após
otimização} pode ser expressa assim:

\begin{quote}
  Para todo programa \(P\) bem tipado,
  \(\llbracket P \rrbracket = \llbracket \mathit{opt}(P) \rrbracket\).
\end{quote}

onde \(\llbracket \cdot \rrbracket\) denota a semântica operacional e
\(\mathit{opt}\) é uma passagem de otimização.

\subsection{Testes de Metamorfose}

\textbf{Metamorphic testing} [Le et al., 2014] aplica transformações
semânticas conhecidas ao programa-fonte e verifica se o compilador
preserva o resultado. Uma relação metamórfica útil é:

\begin{quote}
  Se \(P'\) é obtido de \(P\) por substituição de \(x + 0\) por \(x\),
  então \(P\) e \(P'\) devem produzir o mesmo resultado compilados.
\end{quote}

\section{Verificação Formal de Compiladores}

\subsection{Preservação Semântica}

Um compilador \(\mathcal{C}\) de linguagem-fonte \(\mathcal{L}_s\) para
linguagem-alvo \(\mathcal{L}_t\) é \textbf{correto} se, para todo
programa \(P\) bem tipado e toda entrada \(i\):

\[
P \Downarrow v \implies \mathcal{C}(P) \Downarrow v
\]

onde \(\Downarrow\) denota avaliação segundo a semântica operacional de
cada linguagem. Em palavras: a compilação preserva o comportamento
observável do programa. Esta propriedade é chamada de
\textbf{preservação semântica} (\emph{semantic preservation}).

\subsection{CompCert}

O \textbf{CompCert} [Leroy, 2009] é o primeiro compilador de uso real
com preservação semântica provada formalmente. Ele compila um
subconjunto significativo de C para assembly x86, PowerPC e ARM. Toda
passagem do compilador (análise de fluxo, alocação de
registradores, geração de código) é acompanhada de uma prova
mecânica no assistente de provas \textbf{Coq} de que ela não altera o
comportamento do programa.

O resultado prático é notável: um estudo com CSmith descobriu que
\textbf{nenhum} wrong-code bug existe no CompCert, enquanto GCC e LLVM
apresentaram dezenas.

\subsection{A Estrutura de uma Prova de Compilador}

Cada passagem \(\mathcal{T}\) é acompanhada de uma \textbf{relação de
simulação}:

\[
\infer[\text{Sim}]
  {e \xrightarrow{*} v \implies \mathcal{T}(e) \xrightarrow{*} v}
  {e \xrightarrow{} e' \quad \mathcal{T}(e) \xrightarrow{*} \mathcal{T}(e')}
\]

A composição de simulações para cada passagem produz a preservação
semântica global.

\subsection{Verificação de Passagens Específicas}

Nem sempre é necessário verificar o compilador inteiro. Pesquisas
recentes verificam passagens individuais de alta importância:

\begin{itemize}
  \item \textbf{Alocação de registradores} [Rideau \& Leroy, 2010]:
    uma das passagens mais complexas do GCC foi verificada no Coq.
  \item \textbf{Compilação de JIT} [Carbonneaux et al., 2014]:
    verificação de compiladores just-in-time com garantias de tempo real.
  \item \textbf{Otimizações em LLVM} [Lopes et al., 2015]: a ferramenta
    \textbf{Alive} verifica automaticamente a correção de transformações
    de peephole do LLVM usando SMT solvers.
\end{itemize}

\section{Sistemas de Tipos Avançados}

\subsection{Tipos Dependentes}

Em um sistema de tipos comum, tipos e valores habitam mundos separados:
tipos \emph{classificam} valores, mas valores não podem aparecer em
tipos. \textbf{Tipos dependentes} eliminam essa separação: tipos podem
\emph{depender} de valores.

O exemplo canônico é o vetor de tamanho estático:

\[
\mathsf{Vec} : \mathbb{N} \to \mathsf{Type} \to \mathsf{Type}
\]

\(\mathsf{Vec}\; n\; A\) é o tipo dos vetores com exatamente \(n\)
elementos do tipo \(A\). A função de concatenação tem tipo:

\[
\mathsf{append} : \mathsf{Vec}\; m\; A \to \mathsf{Vec}\; n\; A \to \mathsf{Vec}\; (m + n)\; A
\]

O compilador verifica \emph{em tempo de compilação} que \(m + n\) é o
tamanho correto do resultado. Funções que tentariam acessar um índice
fora do intervalo são rejeitadas pelo verificador de tipos antes de
existir qualquer código gerado.

\textbf{Linguagens com tipos dependentes:} Agda [Norell, 2007], Idris
[Brady, 2013], Coq [Bertot \& Castéran, 2004], Lean 4 [de Moura et
al., 2021].

\textbf{Desafio para compiladores:} inferência de tipos em sistemas
dependentes é em geral indecidível; compiladores exigem anotações
explícitas e resolvem equações em que valores e tipos se misturam.

\subsection{Tipos Lineares}

Um \textbf{tipo linear} garante que um recurso é usado \emph{exatamente
uma vez}: nem descartado sem uso, nem duplicado. A regra de
\textbf{weakening} (descartar uma variável não usada) e a de
\textbf{contraction} (duplicar uma variável) são proibidas para tipos
lineares.

O exemplo motivador é o tratamento de recursos:

\begin{verbatim}
open   : String → File         -- abre arquivo
read   : File   → (String, File)
close  : File   → Unit
\end{verbatim}

Com tipos lineares, o compilador garante que:

\begin{itemize}
  \item Todo arquivo aberto é fechado exatamente uma vez.
  \item Não existe \emph{use after free}.
  \item Não existe \emph{double close}.
\end{itemize}

Se o programador tentar usar \texttt{f} após chamar \texttt{close f},
o verificador de tipos emite um erro em tempo de compilação.

\textbf{Conexão com Rust:} O sistema de \emph{ownership} e
\emph{borrows} do Rust é uma forma de tipos lineares/afins integrada
a uma linguagem de sistemas. O compilador Rust garante ausência de
data races e uso de memória liberada por construção, sem garbage
collector.

\textbf{Base teórica:} Lógica Linear [Girard, 1987];
\emph{Linear Type Theory} [Wadler, 1990];
\emph{Quantitative Type Theory} [Atkey, 2018].

\subsection{Session Types}

\textbf{Session types} tipam canais de comunicação, descrevendo o
\emph{protocolo} que as duas pontas de um canal devem seguir. O tipo
de um canal especifica a sequência de mensagens enviadas e recebidas,
e o compilador verifica que ambas as pontas seguem o protocolo.

O protocolo de um servidor de calculadora pode ser descrito por:

\[
S = \mathsf{!Int}.\; \mathsf{!Int}.\; \mathsf{?Int}.\; \mathsf{end}
\]

Que se lê: ``envie um inteiro, envie outro inteiro, receba um inteiro,
encerre''. O tipo do cliente é o \textbf{dual} \(\bar{S}\):

\[
\bar{S} = \mathsf{?Int}.\; \mathsf{?Int}.\; \mathsf{!Int}.\; \mathsf{end}
\]

O compilador garante que cliente e servidor nunca ficam em
\emph{deadlock} e que não há erros de protocolo (\emph{type mismatches})
em tempo de execução.

\textbf{Linguagens:} Links [Cooper et al., 2006], Scala com
bibliotecas de session types, extensões para Haskell e Rust.
\textbf{Base teórica:} \emph{Session Types} [Honda, 1993];
\emph{A Theory of Communication} [Honda et al., 1998].

\section{Representações Intermediárias}

\subsection{Static Single Assignment (SSA)}

A forma \textbf{SSA} (\emph{Static Single Assignment}) [Cytron et al.,
1991] é uma representação intermediária em que \textbf{cada variável é
atribuída exatamente uma vez}. Programas com atribuições múltiplas são
transformados introduzindo novas versões de cada variável:

\textbf{Antes da SSA:}

\begin{verbatim}
x = 0
x = x + 1
y = x * 2
if b:
    x = 5
else:
    x = 10
z = x + y
\end{verbatim}

\textbf{Após a forma SSA:}

\begin{verbatim}
x0 = 0
x1 = x0 + 1
y0 = x1 * 2
if b:
    x2 = 5
else:
    x3 = 10
x4 = phi(x2, x3)
z0 = x4 + y0
\end{verbatim}

A \textbf{função \(\phi\)} (phi-function) nos pontos de junção de
fluxo de controle seleciona o valor correto da variável dependendo do
caminho percorrido em tempo de execução. Ela não corresponde a uma
instrução de máquina: é uma notação da forma intermediária eliminada
durante a geração de código.

\textbf{Por que SSA facilita otimizações?} Como cada variável é
definida num único ponto, a análise de \emph{use-def chains} se torna
trivial: o único lugar onde \(x_4\) é definido é a phi-function, e o
único lugar onde é usada é \(z_0 = x_4 + y_0\). Isso acelera:

\begin{itemize}
  \item \textbf{Propagação de constantes:} se \(x_0 = 0\) e
    \(x_1 = x_0 + 1\), então \(x_1 = 1\) pode ser propagado
    diretamente.
  \item \textbf{Eliminação de código morto:} variáveis sem uso são
    imediatamente identificáveis.
  \item \textbf{Alocação de registradores:} interferência entre
    variáveis é calculada eficientemente sobre o grafo de interferência
    induzido pela SSA.
\end{itemize}

A SSA é a representação central do \textbf{LLVM IR} e do \textbf{GCC
GIMPLE}.

\subsection{Value State Dependence Graph (VSDG)}

O \textbf{VSDG} [Johnson et al., 2004] é uma representação baseada em
grafos que torna explícitas as dependências entre operações ao invés de
apenas ordená-las sequencialmente. Ao contrário da SSA, que ainda é
essencialmente uma sequência de instruções com anotações, o VSDG
representa o programa como um grafo dirigido acíclico (DAG) com três
tipos de aresta:

\begin{center}
\begin{tabular}{ll}
\toprule
Tipo de aresta & Significado \\
\midrule
\textbf{Valor}   & O resultado de uma operação é entrada de outra \\
\textbf{Estado}  & Uma operação com efeito deve preceder outra (ex.: leituras/escritas) \\
\textbf{Controle}& Uma operação só executa se uma condição for satisfeita \\
\bottomrule
\end{tabular}
\end{center}

\textbf{Exemplo:} o programa

\begin{verbatim}
x = a + b
y = a + b
z = x * y
\end{verbatim}

na SSA geraria dois nós \texttt{add} separados. No VSDG, as duas
computações de \texttt{a + b} são representadas pelo \emph{mesmo} nó,
pois são idênticas. Isso torna a \textbf{eliminação de subexpressões
comuns} (\emph{CSE}) \emph{gratuita}: é uma propriedade estrutural
da representação, não uma passagem separada.

\textbf{Dependências de estado} tornam explícito o sequenciamento de
operações com efeitos colaterais. No VSDG, \texttt{read(x)} e
\texttt{write(x,~v)} têm uma aresta de estado entre si; duas leituras
independentes \emph{não} têm aresta entre si e podem ser reordenadas
livremente pelo compilador.

\textbf{Scheduling:} a partir do VSDG, a ordem de execução é
determinada por um algoritmo de \emph{scheduling} que percorre o grafo
respeitando apenas as dependências reais. Isso habilita reordenações de
instruções que seriam difíceis de detectar na SSA.

\section{Otimização de Código}

\subsection{Passos Clássicos}

As passagens de otimização mais fundamentais são:

\textbf{Constant folding e constant propagation:} expressões cujos
operandos são conhecidos em tempo de compilação são avaliadas pelo
compilador.

\[
\mathit{fold}(3 + 4) = 7 \qquad
\mathit{fold}(\mathit{if}\ \mathsf{true}\ \mathsf{then}\ e_1\ \mathsf{else}\ e_2) = e_1
\]

\textbf{Eliminação de código morto} (\emph{DCE}, \emph{Dead Code
Elimination}): código cujo resultado nunca é observado é removido. Na
forma SSA, uma variável é morta se não tem usos.

\textbf{Inlining:} substituição de uma chamada de função pelo corpo da
função, com os parâmetros formais substituídos pelos argumentos reais.
Elimina o overhead de chamada e expõe outras oportunidades de
otimização.

\textbf{Loop transformations:} rotações, unrolling e vetorização de
laços para explorar o paralelismo de dados (SIMD) dos processadores
modernos.

\subsection{LLVM como Infraestrutura de Pesquisa}

O \textbf{LLVM} [Lattner \& Adve, 2004] tornou-se a principal
plataforma de pesquisa em otimização de compiladores. Sua arquitetura
modular expõe um \emph{pass manager} no qual novas passagens de
otimização podem ser inseridas com poucas linhas de código. O
\textbf{LLVM IR} (uma forma de SSA com tipagem estática) serve de
entrada para dezenas de passagens clássicas e experimentais:

\begin{itemize}
  \item \texttt{mem2reg}: eleva alocações em memória para variáveis
    SSA (a própria transformação SSA).
  \item \texttt{instcombine}: milhares de reescritas algébricas locais.
  \item \texttt{loop-vectorize}: vetorização automática de laços.
  \item \texttt{polly}: otimizações de laços com o modelo poliédrico.
\end{itemize}

\textbf{Ferramenta de validação: Alive2} [Lopes et al., 2021] verifica
automaticamente, via SMT (Z3), se cada transformação do
\texttt{instcombine} preserva a semântica para todas as entradas
possíveis. Bugs encontrados por Alive2 foram corrigidos em versões
oficiais do LLVM.

\subsection{Otimização Guiada por Perfil}

\textbf{Profile-Guided Optimization (PGO)} coleta dados de execuções
reais do programa (frequência de cada bloco básico, alvos de chamadas
indiretas) e os usa para tomar decisões de otimização mais precisas:

\begin{itemize}
  \item Funções frequentemente chamadas são prioritárias para inlining.
  \item Blocos básicos frequentes são colocados em caminhos de execução
    rápidos (\emph{branch layout}).
  \item Tabelas de despacho virtual são ordenadas por frequência.
\end{itemize}

PGO é central em compiladores de produção como o Clang e o JIT do V8
(JavaScript).

\section{Conclusão}

A pesquisa em compiladores é um campo de fronteira que combina
preocupações práticas (desempenho, correção, segurança) com as
ferramentas mais sofisticadas da teoria de linguagens de programação.
As quatro frentes apresentadas neste capítulo são complementares: a
pesquisa em testes descobre bugs existentes em compiladores reais; a
verificação formal elimina categorias inteiras de bugs por construção;
os sistemas de tipos avançados tornam os compiladores capazes de
garantir propriedades mais ricas sobre os programas que aceitam; e as
representações intermediárias e as técnicas de otimização determinam
a qualidade do código gerado.

O estudo formal de compiladores que percorre estas notas de aula, da
análise léxica e sintática à geração de código intermediário, passando
por semântica operacional, sistemas de tipos e inferência de tipos,
é a base necessária para contribuir nessas frentes. Os compiladores
provados corretos, como o CompCert, dependem de semânticas
operacionais rigorosas e de relações de simulação análogas às
estudadas nos capítulos de subtipagem e preservação. Os sistemas de
tipos avançados, como tipos dependentes e lineares, são extensões
diretas dos sistemas de tipos simples e da inferência de Hindley-Milner
apresentados nos capítulos anteriores. A forma SSA e o VSDG são
transformações da representação intermediária estudada na última parte
do texto.

A fronteira mais ativa da área é, atualmente, a interseção entre
verificação formal e otimização: ferramentas como o Alive2 aplicam
raciocínio automático para verificar que transformações de otimização
são corretas para \emph{todas} as entradas possíveis, aproximando o
nível de garantia do CompCert ao mundo das otimizações agressivas do
LLVM. Ao mesmo tempo, a proliferação de linguagens com sistemas de
tipos mais expressivos (Rust com \emph{ownership}, Lean 4 com tipos
dependentes, Haskell com tipos lineares opcionais) coloca novos
desafios para a teoria e a implementação de compiladores, desafios que
as ferramentas formais desenvolvidas nos capítulos anteriores estão
bem posicionadas para enfrentar.

\section{Referências}

\subsection{Teste de Compiladores}

\begin{itemize}
  \item Yang, X., Chen, Y., Eide, E., \& Regehr, J. (2011).
    \textbf{Finding and understanding bugs in C compilers.}
    \emph{PLDI 2011}.
  \item Claessen, K., \& Hughes, J. (2000).
    \textbf{QuickCheck: a lightweight tool for random testing of
    Haskell programs.} \emph{ICFP 2000}.
  \item Le, V., Afshari, M., \& Su, Z. (2014).
    \textbf{Compiler validation via equivalence modulo inputs.}
    \emph{PLDI 2014}.
\end{itemize}

\subsection{Verificação Formal}

\begin{itemize}
  \item Leroy, X. (2009).
    \textbf{Formal verification of a realistic compiler.}
    \emph{CACM 52(7)}.
  \item Lopes, N. P., Menendez, D., Nagarakatte, S., \& Regehr, J.
    (2015). \textbf{Provably correct peephole optimizations with Alive.}
    \emph{PLDI 2015}.
  \item Rideau, L., \& Leroy, X. (2010).
    \textbf{Validating register allocation and spilling.} \emph{CC 2010}.
\end{itemize}

\subsection{Sistemas de Tipos Avançados}

\begin{itemize}
  \item Pierce, B. C. (2002).
    \textbf{Types and Programming Languages.} MIT Press.
    \emph{(Capítulos 23--32 cobrem sistemas avançados de tipos.)}
  \item Brady, E. (2013).
    \textbf{Idris, a general purpose dependently typed programming
    language.} \emph{Journal of Functional Programming 23(5)}.
  \item Wadler, P. (1990).
    \textbf{Linear types can change the world!}
    \emph{IFIP TC 2 Working Conference on Programming Concepts and
    Methods}.
  \item Honda, K., Vasconcelos, V. T., \& Kubo, M. (1998).
    \textbf{Language primitives and type discipline for structured
    communication-based programming.} \emph{ESOP 1998}.
\end{itemize}

\subsection{Formas Intermediárias e Otimização}

\begin{itemize}
  \item Cytron, R., Ferrante, J., Rosen, B. K., Wegman, M. N., \&
    Zadeck, F. K. (1991). \textbf{Efficiently computing static single
    assignment form and the control dependence graph.}
    \emph{TOPLAS 13(4)}.
  \item Johnson, N., \& Mycroft, A. (2004).
    \textbf{Combined code motion and register allocation using the
    Value State Dependence Graph.} \emph{CC 2004}.
  \item Lattner, C., \& Adve, V. (2004).
    \textbf{LLVM: a compilation framework for lifelong program analysis
    and transformation.} \emph{CGO 2004}.
  \item Lopes, N. P., Lee, J., Hur, C.-K., Liu, Z., \& Regehr, J.
    (2021). \textbf{Alive2: bounded translation validation for LLVM.}
    \emph{PLDI 2021}.
\end{itemize}

\subsection{Leituras Complementares}

\begin{itemize}
  \item Appel, A. W. (1998).
    \textbf{Modern Compiler Implementation in ML.}
    Cambridge University Press.
  \item Cooper, K., \& Torczon, L. (2011).
    \textbf{Engineering a Compiler} (2ª ed.). Morgan Kaufmann.
  \item Winskel, G. (1993).
    \textbf{The Formal Semantics of Programming Languages.} MIT Press.
\end{itemize}
